\documentclass[a4paper,fleqn]{cas-sc}

\usepackage[authoryear,longnamesfirst]{natbib}
\usepackage{amsmath}
\usepackage{amssymb}
\usepackage{amsthm}
\usepackage{booktabs}
\usepackage{array}
\usepackage{subcaption}
\usepackage{siunitx}
\usepackage{longtable}
\usepackage{xcolor}
\usepackage{lineno}
\usepackage[ruled,vlined]{algorithm2e}

% ============================================================
% Theorem environments
% ============================================================
\newtheorem{theorem}{Theorem}
\newtheorem{proposition}{Proposition}
\newtheorem{corollary}{Corollary}

% ============================================================
% Compact float/table controls
% ============================================================
\renewcommand{\topfraction}{0.90}
\renewcommand{\bottomfraction}{0.80}
\renewcommand{\textfraction}{0.08}
\renewcommand{\floatpagefraction}{0.75}
\setcounter{topnumber}{2}
\setcounter{bottomnumber}{1}
\setcounter{totalnumber}{4}

\newcommand{\figsmall}{0.58\textwidth}
\newcommand{\figmedium}{0.68\textwidth}
\newcommand{\figlarge}{0.78\textwidth}
\newcommand{\figwide}{0.84\textwidth}
\newcommand{\figappendix}{0.78\textwidth}

\newcommand{\tablefont}{\small}
\newcommand{\tighttables}{%
  \setlength{\tabcolsep}{4.5pt}%
  \renewcommand{\arraystretch}{1.08}%
}

\newcommand{\tinytable}{\scriptsize}
\newcommand{\verytighttables}{%
  \setlength{\tabcolsep}{3.5pt}%
  \renewcommand{\arraystretch}{1.05}%
}

% ============================================================
% Macros
% ============================================================
\newcommand{\x}{\mathbf{x}}
\newcommand{\z}{\mathbf{z}}
\newcommand{\h}{\mathbf{h}}
\newcommand{\uvar}{\mathbf{u}}
\newcommand{\y}{\mathbf{y}}
\newcommand{\R}{\mathbb{R}}
\newcommand{\Hspace}{\mathcal{H}}
\newcommand{\Zspace}{\mathcal{Z}}
\newcommand{\Sset}{\mathcal{S}}
\newcommand{\Lrec}{\mathcal{L}_{\mathrm{rec}}}
\newcommand{\Lmono}{\mathcal{L}_{\mathrm{mono}}}
\newcommand{\Lsmooth}{\mathcal{L}_{\mathrm{smooth}}}

% Artifact roots
\newcommand{\artifactroot}{paper_artifacts/acml_results_bundle}
\newcommand{\mainfig}{\artifactroot/main/figures}
\newcommand{\maintab}{\artifactroot/main/tables}
\newcommand{\appfig}{\artifactroot/appendix/figures}
\newcommand{\apptab}{\artifactroot/appendix/tables}
\newcommand{\resultsfig}{results/figures}
\newcommand{\resultstab}{results/tables}

\begin{document}

\let\WriteBookmarks\relax
\def\floatpagepagefraction{1}
\def\textpagefraction{.001}

% ============================================================
% Front matter
% ============================================================
\shorttitle{Learning Bounded Latent Degradation Dynamics}
\shortauthors{A. Goel}

\title[mode=title]{Learning Bounded Latent Degradation Dynamics for stable rollout and Remaining Useful Life Prediction}

\author[1]{Arpit Goel}
\cormark[1]
\ead{arpit@thebrownkid.in}

\affiliation[1]{organization={TwinSim Labs}}

\cortext[1]{Corresponding author}

\begin{abstract}
Long-horizon multivariate time-series forecasting is often unstable when models
are recursively rolled out in observation space. This problem is especially
pronounced in degradation modelling, where sensor trajectories are noisy,
regime-dependent, and only partially informative about latent health. We study
whether low-dimensional bounded latent representations can act as dynamical
regularizers for stable rollout. The proposed algorithm identifies operating
regimes from control variables, normalizes sensor measurements within regimes,
denoises the resulting trajectories, selects degradation-sensitive channels,
and trains a bounded latent autoencoder with monotonicity and smoothness
regularization on a primary health coordinate. We provide a theoretical
analysis separating sensor-space closed-loop divergence, bounded latent
rollout, and smoothness-controlled forecastability, and we extend it in two
directions: an error bound for damped and windowed extrapolation, and anchored
decoding, which cancels the reconstruction offset that a decoded forecast
carries at its origin. Experiments on NASA C-MAPSS show that
higher-dimensional latent spaces can improve reconstruction while weakening
rollout stability. Ablations further show that finite-horizon decoded
boundedness is not unique to the bounded latent model: an unbounded
autoencoder may also remain numerically bounded under constant-velocity
rollout because of small latent velocities and decoder saturation. The main
empirical advantage of the proposed representation is therefore forecastability:
lower latent curvature, higher forecast skill, and useful downstream
remaining-useful-life prediction. On the PHM~2010 CNC milling benchmark both
mechanisms transfer outside turbofan degradation: the sensor-space recursion
leaves the boundedness criterion under a long free run while every bounded
latent rollout stays controlled, and anchored decoding removes the
short-horizon penalty that plain decoding incurs.
\end{abstract}

\begin{keywords}
stable forecasting \sep latent representation learning \sep multivariate time series \sep autoencoders \sep degradation modelling \sep remaining useful life
\end{keywords}

\maketitle

\section{Introduction}
% ============================================================

Many multivariate time-series applications require models that remain useful
when rolled forward beyond a one-step training objective. In degradation
modelling and prognostics, this requirement is especially important: forecasts
are used not only to interpolate observed measurements, but also to reason
about future system health. A model that reconstructs historical trajectories
accurately may still produce implausible or unstable behavior when recursively
rolled out.

Raw sensor measurements mix several sources of variation. Measurements change
with operating conditions, sensor noise, unit-specific behavior, and latent
degradation. Directly forecasting in observation space can therefore lead to
unstable dynamics. Linear autoregressive models may become non-contractive
when their closed-loop transition operators have spectral radius greater than
one \citep{hamilton1994time,lutkepohl2005new}, while high-capacity sequence
models may achieve good short-horizon accuracy without guaranteeing stable
free-run behavior, a difficulty long documented for recurrent architectures
trained via backpropagation through time \citep{pascanu2013difficulty}.

This paper studies an algorithmic alternative. Instead of forecasting raw
sensors directly, we learn a compact latent state intended to represent
degradation after removing operating-regime effects. The latent state is
bounded by construction and is regularized so that its primary coordinate
evolves smoothly and approximately monotonically over the life of each unit.
Forecasting is then performed by local constant-velocity rollout in latent
space, followed by decoding back to the sensor representation.

The key thesis is that stable long-horizon forecasting requires more than
accurate reconstruction. Bounded latent geometry constrains the admissible state
space, while smooth low-curvature latent dynamics make motion within that
space forecastable. These are distinct mechanisms: boundedness alone is not
sufficient for useful forecasting, and short-horizon accuracy does not ensure
stable free-run behavior.

The main contributions are:
\begin{itemize}
    \item We propose a regime-aware latent representation learning algorithm for
    multivariate degradation time series, combining condition normalization,
    denoising, dynamic sensor selection, and bounded latent autoencoding.

    \item We provide a theoretical analysis separating sensor-space
    closed-loop divergence, bounded latent rollout, and smoothness-controlled
    latent forecastability.

    \item We show empirically that boundedness alone is insufficient: an
    unbounded latent autoencoder can remain finite over practical horizons,
    but exhibits higher latent curvature and weaker forecast skill.

    \item We demonstrate that increasing latent dimension improves
    reconstruction but can weaken rollout stability, supporting
    low-dimensional latent manifolds as dynamical regularizers.

    \item We compare against VAR, VAR(2), GRU, LSTM, and TCN baselines, showing
    that closed-loop sensor-space and neural sequence forecasters can exhibit
    inconsistent free-run stability.

    \item We introduce \emph{anchored decoding}, a reformulation of the decoded
    forecast that cancels the origin reconstruction offset. We prove its error
    obeys the same curvature bound as the underlying rollout up to the decoder
    Lipschitz constant, plus a reconstruction-residual term, and we give the
    condition on that residual under which anchoring is an improvement.
    Empirically it removes the short-horizon penalty that plain decoding incurs,
    changing horizon-one skill on CNC tool wear from $-263$ to $-0.73$.
\end{itemize}

% ============================================================
\section{Related Work}
\label{sec:related}
% ============================================================

\subsection{RUL Prediction and Health Indicators}

Remaining useful life (RUL) estimation is a central task in prognostics and
health management, where the goal is to infer the remaining time before a
system reaches end of useful life from condition-monitoring data
\citep{si2011remaining}. NASA C-MAPSS is a standard benchmark for this setting:
it provides run-to-failure multivariate time series, operating settings,
sensor noise, official train/test splits, and test RUL labels
\citep{nasa_cmapss,saxena2008turbofan,saxena2008damage}. The four subsets FD001--FD004 vary in
operating-condition complexity and fault-mode multiplicity.

Classical health-indicator construction emphasizes properties such as
monotonicity, trendability, and prognosability
\citep{coble2009identifying}. These criteria motivate representations that move
consistently with degradation and can be extrapolated toward failure. More
recent work explicitly builds monotonicity and other domain constraints into
the learning objective for data-driven health indicators, closer in spirit to
the regularization used here \citep{tefera2025constraint}. Our work
adopts this health-indicator perspective but learns the health state as a
latent coordinate rather than hand-engineering it from individual sensors.

\subsection{Deep Learning for Multivariate Prognostics}

Deep learning methods have been widely used for RUL prediction from
multivariate sensor windows. Convolutional models learn local temporal features
\citep{babu2016deep,li2018deepcnn}, while recurrent architectures such as LSTMs
and GRUs model sequential dependencies
\citep{zheng2017lstm,hochreiter1997lstm,cho2014gru}. Temporal convolutional
networks provide a strong generic sequence-modelling baseline with causal
dilated convolutions and residual connections \citep{bai2018tcn}. Attention
mechanisms have also been layered onto recurrent and convolutional backbones
for RUL estimation, including feature-attention LSTMs
\citep{liu2021featureattention}, attention-augmented LSTMs for rotating
machinery \citep{zhang2020attentionlstm}, and bidirectional GRUs with temporal
self-attention \citep{zhang2022bigru}. Recent C-MAPSS RUL models include
CNN-LSTM-Attention \citep{deng2024cnnlstmattention}
and Transformer--TCN--Self-Attention feature fusion \citep{li2025ttsnet}.

Beyond turbofan data, the PHM~2010 high-speed milling benchmark \citep{phm2010} is
a standard testbed for tool-wear prognostics from force, vibration and
acoustic-emission channels, and is commonly addressed with deep sequence models
trained on its labelled cutters \citep{zhang2023toolwear}. We use it in
Section~\ref{sec:milling} as a degradation process whose sensor dynamics differ
substantially from C-MAPSS.

However, supervised RUL accuracy and stable free-run forecasting are different
objectives. A sequence model can be accurate over a short horizon or under
teacher forcing while becoming unstable when recursively rolled forward. In
this paper, GRU, LSTM, and TCN models are treated as forecasting baselines for
testing free-run stability rather than as direct replacements for the proposed
latent representation.

\subsection{Autoencoders and Latent Dynamics}

Autoencoders are commonly used to learn compact representations of
high-dimensional signals. Denoising autoencoders learn features robust to input
corruption \citep{vincent2008dae}, and variational autoencoders introduce
probabilistic latent variables via amortized stochastic inference
\citep{kingma2014autoencoding,rezende2014stochastic}. In prognostics,
autoencoder representations are often used as features for downstream health
or RUL prediction.

Latent dynamics methods similarly separate high-dimensional observations from
lower-dimensional evolution. This perspective traces back to classical operator
theory, in which a nonlinear dynamical system is represented through the
linear evolution of observables in an infinite-dimensional function space
\citep{koopman1931hamiltonian}. Koopman autoencoders operationalize this idea
with finite-dimensional learned embeddings in which nonlinear dynamics are
approximated by linear evolution
\citep{lusch2018deep,erichson2020forecasting}. Latent Space Dynamics
Identification combines autoencoder compression with learned latent dynamics
for reduced-order rollout \citep{fries2022lasdi}, with subsequent work adding
Gaussian-process uncertainty over the latent dynamics
\citep{bonneville2024gplasdi}. Our work shares the broad
principle of forecasting in latent space, but focuses on how bounded
low-dimensional geometry and latent smoothness affect free-run stability and
forecastability in degradation time series.

\subsection{Position of This Work}

The contribution of this paper is not a new autoencoder architecture alone.
Rather, we study bounded latent manifolds as dynamical regularizers. The
experiments explicitly separate reconstruction, boundedness, curvature,
forecast skill, and downstream RUL. This distinction is important because a
model may reconstruct well but forecast poorly, and a model may remain
numerically bounded over finite horizons without learning a useful degradation
trajectory.

Koopman autoencoders and LaSDI share the premise that forecasting should happen
in a reduced space, but they target a different data setting. Both fit a latent
evolution operator to full-state snapshots of a system whose dynamics are
observed consistently: LaSDI to solution snapshots of a parameterized
simulation, Koopman autoencoders to repeated observations of a single
autonomous system. C-MAPSS supplies neither. Each unit is a separate,
right-censored trajectory with unknown initial wear, sampled under discrete
operating regimes and observed only through partial noisy sensors, with no
shared parameterization across units and no governing equations. These methods
are therefore not baselines that can be transferred to this setting, and we do
not present them as such.

% ============================================================
\section{Method}
\label{sec:method}
% ============================================================

\subsection{Problem Setup}

Let $\x_{i,t}\in\R^p$ denote the sensor vector for unit $i$ at cycle $t$, and
let $\uvar_{i,t}\in\R^q$ denote the corresponding operating variables. The goal
is to learn a low-dimensional latent state $\h_{i,t}\in\R^K$, with $K\ll p$,
that captures degradation-relevant dynamics and supports stable long-horizon
rollout.

The workflow is summarized in Fig.~\ref{fig:pipeline}. It first removes
nuisance variation caused by operating conditions and static sensors, then
learns a compact latent state whose geometry is suitable for free-run
forecasting.

\begin{figure}[t]
    \centering
    \includegraphics[width=\figlarge,height=0.38\textheight,keepaspectratio]{\mainfig/fig_pipeline.png}
    \caption{Regime-aware bounded latent manifold learning.}
    \label{fig:pipeline}
\end{figure}

\begin{algorithm}[t]
\caption{Bounded Latent Degradation Manifold Learning}
\label{alg:manifold}
\KwIn{Trajectories $\{\x_{i,t},\uvar_{i,t}\}$ and latent dimension $K$}
\KwOut{Bounded latent state $\h_{i,t}$, rollout $\hat{\h}_{i,t+\tau}$, and RUL features}
Fit operating-regime map $r_{i,t}=C(\uvar_{i,t})$ using control variables only\;
Estimate regime-specific statistics and compute normalized residuals\;
Denoise residual trajectories using a rolling median\;
Select dynamic sensors using degradation-trend scores\;
Train bounded autoencoder
$\h_{i,t}=E_\theta(\y_{i,t})\in[0,1]^K$,
$\hat{\y}_{i,t}=D_\phi(\h_{i,t})$\;
Penalize reconstruction error, negative increments of $h_0$, and roughness of $h_0$\;
Estimate local latent velocity and roll out
$\hat{\h}_{t+\tau}=\Pi_{[0,1]^K}(\h_t+\tau\hat{\mathbf v}_t)$\;
Decode $\hat{\h}_{t+\tau}$ or use $[\h,\dot{\h}]$ for RUL prediction\;
\end{algorithm}

\subsection{Regime Identification and Normalization}

Operating conditions can dominate raw sensor variation. We therefore identify
discrete regimes using only the operating variables. A clustering function $C$
assigns each observation to a regime, $r_{i,t}=C(\uvar_{i,t})$. In the
implementation, $C$ is a K-means model fit on the three operating settings.
Sensors are not used to form regimes, preventing degradation-sensitive
measurements from influencing the operating-condition partition.

For each regime $r$ and sensor $j$, we estimate a regime-specific mean
$\mu_{r,j}$ and scale $\sigma_{r,j}$. The normalized residual is
\[
    \tilde{x}_{i,t,j}
    =
    \frac{x_{i,t,j}-\mu_{r_{i,t},j}}
    {\sigma_{r_{i,t},j}+\epsilon}.
\]
This removes operating-regime offsets and yields trajectories more likely to
reflect degradation.

\subsection{Denoising and Dynamic Sensor Selection}

The normalized trajectories are denoised independently for each unit using a
rolling median,
\[
    \bar{x}_{i,t,j}
    =
    \operatorname{median}_{\tau\in W(t)}
    \tilde{x}_{i,\tau,j}.
\]
This suppresses local spikes while preserving slowly varying degradation
trends.

Not every sensor channel is degradation-informative. For each denoised sensor
$j$, we compute
\[
    S_j
    =
    \frac{1}{N}
    \sum_i
    \left|
    \operatorname{corr}(\bar{x}_{i,:,j},t)
    \right|.
\]
The selected dynamic sensor set is $\Sset=\{j:S_j\geq\theta\}$, and the
autoencoder input is $\y_{i,t}=\bar{\x}_{i,t,\Sset}$. Thus the autoencoder is
not asked to spend capacity reconstructing static or purely
operating-condition channels.

\subsection{Bounded Latent Autoencoder}

The encoder maps the selected denoised sensor vector to a bounded latent state:
\[
    \h_{i,t}=E_\theta(\y_{i,t})\in[0,1]^K,
    \qquad
    \hat{\y}_{i,t}=D_\phi(\h_{i,t}).
\]
The objective is
\[
    \mathcal{L}
    =
    \Lrec+\lambda_m\Lmono+\lambda_s\Lsmooth,
    \qquad
    \Lrec=\sum_{i,t}\|\hat{\y}_{i,t}-\y_{i,t}\|_2^2.
\]
The monotonicity and smoothness penalties are applied to the primary latent
coordinate $h_0$:
\[
    \Lmono
    =
    \sum_{i,t}
    \max(0,-\Delta h_{i,t,0}),
    \qquad
    \Lsmooth
    =
    \sum_{i,t}
    (\Delta h_{i,t,0})^2,
\]
where $\Delta h_{i,t,0}=h_{i,t+1,0}-h_{i,t,0}$. After training, $h_0$ is
oriented so that it increases with cycle index. The smoothness penalty follows
the broader principle, used in physics-informed and constraint-regularized
learning, of penalizing a differential-operator norm of the model output to
encourage solutions that are physically or dynamically plausible rather than
merely low-error on observed points \citep{raissi2019physics}.

\subsection{Latent Rollout and RUL Features}

Forecasting is performed in latent space. A local velocity
$\hat{\mathbf{v}}_t$ is estimated over a trailing window, and future latent
states are predicted as
\[
    \hat{\h}_{t+\tau}
    =
    \Pi_{\Hspace}
    \left(
    \h_t+\tau\hat{\mathbf{v}}_t
    \right),
    \qquad
    \Hspace=[0,1]^K .
\]
The decoded forecast is $\hat{\y}_{t+\tau}=D_\phi(\hat{\h}_{t+\tau})$.
All bounded-rollout results use full-vector projection, matching the
compact-domain assumption in Theorem~\ref{thm:bounded-rollout}. A legacy
$h_0$-only clipping rule was also tested and produced materially similar
metrics.

Constant-velocity rollout is the simplest motion model that isolates the effect
of latent geometry, and it is what all C-MAPSS results in
Section~\ref{sec:results} use. Where degradation does not progress at a steady
rate, we also consider a learned bounded residual alternative,
\[
    \hat{\h}_{t+1}
    =
    \Pi_{\Hspace}
    \bigl(
    \h_t+\gamma\tanh g_\psi([\h_t,\mathbf{c}_t])
    \bigr),
\]
where $g_\psi$ is a small residual MLP trained on the frozen latent
trajectories with a combined one-step and multi-step free-run objective, and
the step size $\gamma$ caps the per-step displacement. The optional input
$\mathbf{c}_t$ carries exogenous variables that are known in advance for every
future step, so the head remains usable at arbitrary horizons without observing
future sensors; where no such variables exist, $g_\psi$ takes $\h_t$ alone.
Because the projection $\Pi_{\Hspace}$ is unchanged,
Theorem~\ref{thm:bounded-rollout} applies verbatim for any $g_\psi$: the
learned variant changes only which motion within $\Hspace$ is predicted, never
whether the rollout stays in $\Hspace$. This head is exercised in the step-size
sweep of Appendix~\ref{app:sensitivity}. It is not of the linear-extrapolation form
analysed in Proposition~\ref{prop:general-extrap}, so unlike the rules below it
carries the boundedness guarantee without a curvature-based error bound, and the
reported rollouts do not use it.

The reported rollouts share a single form. The velocity $\hat{\mathbf{v}}$ is
estimated by ordinary least squares over a trailing window, and the forecast is
\[
    \hat{\h}_{c+\tau}
    =
    \Pi_{\Hspace}\bigl(\h_c+\hat{\mathbf{v}}\,S_\tau\bigr),
\]
where the weights $S_\tau$ select the rule. Taking $S_\tau=\tau$ gives constant
velocity. Taking $S_\tau=(1-\phi^\tau)/(1-\phi)$ damps the extrapolation
geometrically toward a frozen state, in the manner of damped-trend exponential
smoothing \citep{gardner1985forecasting}; $\phi{=}1$ recovers constant velocity and
$\phi{=}0$ freezes the state after a single step. A Holt damped-trend recursion,
which smooths a level and a trend rather than fitting one windowed slope, is the
third member of the family. The damping and smoothing constants are selected on
training trajectories only. Every rule of this form obeys both
Theorem~\ref{thm:bounded-rollout} and Proposition~\ref{prop:general-extrap}.

For RUL prediction, the downstream regressor uses all latent coordinates and
their local velocities:
\[
    [h_0,\ldots,h_{K-1},v_0,\ldots,v_{K-1}].
\]
This avoids confounding the latent-dimension study by discarding coordinates
beyond $h_1$.

\subsection{Ablation Variants}

The comparative study evaluates the full model and controlled variants: no
regime normalization, no monotonicity penalty, no smoothness penalty, no
monotonicity or smoothness, bounded latent autoencoder without penalties, and
an unbounded latent autoencoder. These variants isolate which parts of the
algorithm affect reconstruction, finite-horizon boundedness, latent curvature,
forecast skill, and downstream RUL.

% ============================================================
\section{Theory}
\label{sec:theory}
% ============================================================

This section separates three mechanisms: sensor-space closed-loop instability,
bounded latent rollout, and smoothness-controlled forecastability. The sensor-space
analysis follows standard linear time-series and system-identification theory
\citep{hamilton1994time,lutkepohl2005new,ljung1999system}, while the latent-space
boundedness guarantee relies on elementary point-set topology
\citep{munkres2000topology}.

\subsection{Sensor-Space Closed-Loop Instability}

Consider a vector autoregression in standardized sensor space:
\[
    \z_{t+1}=\sum_{\ell=1}^{p}A_\ell \z_{t+1-\ell}+b .
\]
With stacked state $\xi_t=[\z_t^\top,\z_{t-1}^\top,\ldots]^\top$, the
closed-loop rollout is $\xi_{t+1}=C\xi_t+\tilde{b}$, where $C$ is the
companion matrix.

\begin{theorem}[Non-contractive linear rollout]
\label{thm:linear-growth}
If $\rho(C)>1$, where $\rho(C)$ is the spectral radius of the companion matrix,
then the closed-loop rollout is non-contractive. For generic initial errors
with nonzero projection onto the dominant eigenspace, free-running error grows
geometrically at a rate governed by $\rho(C)$.
\end{theorem}

\begin{proof}
Let $e_t$ be the stacked forecast error. Then $e_{t+1}=Ce_t+r_t$, where $r_t$
is the one-step residual. Iterating gives
\[
    e_t=C^te_0+\sum_{s=0}^{t-1}C^sr_{t-1-s}.
\]
By Gelfand's formula, $\lim_{t\to\infty}\|C^t\|^{1/t}=\rho(C)$. If
$\rho(C)>1$, the homogeneous component grows geometrically for generic initial
errors with nonzero dominant-eigenspace projection.
\end{proof}

\subsection{Bounded Latent Rollout}

Let the latent forecast be
$\hat{\h}_{c+\tau}=\Pi_{\Hspace}(\h_c+\tau\nu_c)$, where
$\Hspace\subset\R^K$ is compact.

\begin{theorem}[Bounded latent rollout]
\label{thm:bounded-rollout}
If $\hat{\h}_{c+\tau}\in\Hspace$ for all $\tau\geq0$ and the decoder
$D_\phi$ is continuous on $\Hspace$, then $D_\phi(\hat{\h}_{c+\tau})$ is
uniformly bounded for all horizons.
\end{theorem}

\begin{proof}
Because $\Hspace$ is compact and $D_\phi$ is continuous, the image
$D_\phi(\Hspace)$ is compact \citep{munkres2000topology}. Compact subsets of finite-dimensional Euclidean
space are bounded, so the decoded rollout is uniformly bounded.
\end{proof}

\begin{corollary}[Boundedness is architectural]
The boundedness guarantee depends on projection into a compact latent domain
and decoder continuity. It does not depend on the monotonicity or smoothness
penalty weights.
\end{corollary}

This guarantee is a formal statement about the bounded latent model. It does
not imply that every unbounded latent model must diverge over a finite horizon:
finite-horizon decoded boundedness can also occur because of small
extrapolation velocities or saturating nonlinearities.

\subsection{Smoothness and Forecastability}

\begin{theorem}[Constant-velocity forecast error]
\label{thm:smooth-forecast}
Let $h_t$ be a scalar latent coordinate with bounded second difference
$|\Delta^2h_t|=|h_{t+1}-2h_t+h_{t-1}|\leq \kappa$. For the
constant-velocity forecast
$\hat{h}_{c+\tau}=h_c+\tau(h_c-h_{c-1})$, the error satisfies
\[
    |h_{c+\tau}-\hat{h}_{c+\tau}|
    \leq
    \frac{\kappa}{2}\tau(\tau+1).
\]
\end{theorem}

\begin{proof}
Let $\delta_s=h_{c+s}-h_{c+s-1}$. Then
\[
    h_{c+\tau}-\hat{h}_{c+\tau}
    =
    \sum_{s=1}^{\tau}(\delta_s-\delta_0).
\]
Since $|\Delta^2h_t|\leq\kappa$, we have
$|\delta_s-\delta_0|\leq s\kappa$. Summing over $s=1,\ldots,\tau$ gives the
bound.
\end{proof}

Thus smoothness does not create boundedness; it reduces the curvature constant
that controls constant-velocity forecast error. The same qualitative point ---
that a bounded second difference controls extrapolation error over a finite
horizon --- underlies classical forecasting practice for stochastic processes
more generally \citep{brockwell2016introduction}.

Theorem~\ref{thm:smooth-forecast} is stated for the textbook
constant-velocity rule, in which the velocity is a single backward difference
and the extrapolation weight is $\tau$ itself. The rollouts used in practice
estimate the velocity over a trailing window and may damp it, so we state the
bound they actually satisfy.

\begin{proposition}[Damped and windowed extrapolation]
\label{prop:general-extrap}
Let $h_t$ satisfy $|\Delta^2h_t|\le\kappa$ and write
$\delta_0=h_c-h_{c-1}$. Consider any rollout of the form
$\hat{h}_{c+\tau}=h_c+v\,S_\tau$, where $v$ is an arbitrary velocity
estimate and $\{S_\tau\}_{\tau\ge1}$ are deterministic weights. Then
\[
    |h_{c+\tau}-\hat{h}_{c+\tau}|
    \le
    \frac{\kappa}{2}\tau(\tau+1)
    +
    |\tau\delta_0-v\,S_\tau| .
\]
\end{proposition}

\begin{proof}
With $\delta_s=h_{c+s}-h_{c+s-1}$ we have
$h_{c+\tau}-h_c=\sum_{s=1}^{\tau}\delta_s
=\tau\delta_0+\sum_{s=1}^{\tau}(\delta_s-\delta_0)$, and
$|\delta_s-\delta_0|\le s\kappa$ as in
Theorem~\ref{thm:smooth-forecast}, so the second sum is at most
$\frac{\kappa}{2}\tau(\tau+1)$ in absolute value. Subtracting
$\hat{h}_{c+\tau}-h_c=v\,S_\tau$ and applying the triangle inequality gives the
stated bound.
\end{proof}

Setting $v=\delta_0$ and $S_\tau=\tau$ makes the second term vanish and recovers
Theorem~\ref{thm:smooth-forecast}. The term is therefore the price of departing
from the textbook rule: a windowed velocity estimate pays $|\tau(\delta_0-v)|$
in exchange for lower variance in $v$, and geometric damping in the manner of
damped-trend exponential smoothing \citep{gardner1985forecasting}
$S_\tau=(1-\phi^\tau)/(1-\phi)$ pays $|v||\tau-S_\tau|$ in exchange for a
rollout that cannot run away. All rollout rules reported in this paper are of
this form, so all of them inherit both this error bound and the boundedness
guarantee of Theorem~\ref{thm:bounded-rollout}.

\subsection{Latent Dimension as Dynamical Regularization}

\begin{proposition}[Capacity-stability tradeoff]
\label{prop:k-regularization}
Increasing latent dimension $K$ weakly increases reconstruction capacity,
because lower-dimensional representations can be embedded in higher-dimensional
ones. At the same time, it increases the number of independent latent
directions available for rollout. Therefore, larger $K$ can improve
reconstruction while weakening dynamical regularization.
\end{proposition}

The $K$-aware dimension sweep tests this proposition using all latent
coordinates and velocities downstream.

\subsection{Anchored Decoding}
\label{sec:anchored}

The forecasts above decode the rolled-out latent directly, $D_\phi(\hat{\h}_{c+\tau})$,
and therefore inherit a fixed reconstruction offset at the forecast origin: even
a perfect latent forecast decodes to $\y_c$ with error $\y_c-D_\phi(\h_c)$. We
remove this offset by anchoring the decoded \emph{increment} to the last
observation.

\begin{theorem}[Anchored decoding forecast error]
\label{thm:anchored}
Let $\hat{\h}_{c+\tau}\in\Hspace$ be a latent rollout and define the anchored
decoded forecast
\[
  \hat{\y}_{c+\tau} = \y_c + \big(D_\phi(\hat{\h}_{c+\tau}) - D_\phi(\h_c)\big),
\]
with reconstruction residual $\mathbf{r}_t=\y_t-D_\phi(\h_t)$. If $D_\phi$ is
$L$-Lipschitz on $\Hspace$ and each latent coordinate satisfies
$|\Delta^2 h^{(j)}_t|\le\kappa$, then
\[
  \|\y_{c+\tau}-\hat{\y}_{c+\tau}\|
  \le L\sqrt{K}\,\tfrac{\kappa}{2}\,\tau(\tau+1) + \|\mathbf{r}_{c+\tau}-\mathbf{r}_c\|.
\]
If in addition $\|\mathbf{r}_t\|\le\rho$ for all $t$, the last term is at most $2\rho$.
\end{theorem}

\begin{proof}
Write $\y_t=D_\phi(\h_t)+\mathbf{r}_t$. The $D_\phi(\h_c)$ terms cancel, giving
\[
  \y_{c+\tau}-\hat{\y}_{c+\tau}
  = \big(D_\phi(\h_{c+\tau})-D_\phi(\hat{\h}_{c+\tau})\big)
  + (\mathbf{r}_{c+\tau}-\mathbf{r}_c).
\]
By $L$-Lipschitz continuity,
$\|D_\phi(\h_{c+\tau})-D_\phi(\hat{\h}_{c+\tau})\|\le L\|\h_{c+\tau}-\hat{\h}_{c+\tau}\|$,
and applying Proposition~\ref{prop:general-extrap} coordinatewise gives
$\|\h_{c+\tau}-\hat{\h}_{c+\tau}\|\le\sqrt{K}\,\tfrac{\kappa}{2}\tau(\tau+1)$. The
triangle inequality with $\|\mathbf{r}_t\|\le\rho$ bounds the residual term by $2\rho$.
\end{proof}

\begin{corollary}[Anchoring cancels the origin offset]
\label{cor:anchor}
Plain decoding $\hat{\y}^{\mathrm{pl}}_{c+\tau}=D_\phi(\hat{\h}_{c+\tau})$ satisfies
$\|\y_{c+\tau}-\hat{\y}^{\mathrm{pl}}_{c+\tau}\|\le
L\sqrt{K}\,\tfrac{\kappa}{2}\tau(\tau+1)+\|\mathbf{r}_{c+\tau}\|$. Anchoring replaces
$\|\mathbf{r}_{c+\tau}\|$ by $\|\mathbf{r}_{c+\tau}-\mathbf{r}_c\|$, which vanishes at
$\tau=0$ and is small whenever the reconstruction residual is temporally
persistent. The leading curvature term is unchanged, so anchoring retains the
bounded-rollout guarantees of Theorems~\ref{thm:bounded-rollout}
and~\ref{thm:smooth-forecast} while eliminating the constant reconstruction
offset.
\end{corollary}

Whether that replacement is an improvement is a property of the residual
process, not a universal fact, and it is measurable before any forecast is
made.

\begin{proposition}[When anchoring helps]
\label{prop:anchor-gain}
Suppose the reconstruction residual is weakly stationary with
$\mathbb{E}\|\mathbf{r}_t\|^2=\sigma^2$ and lag-$\tau$ autocorrelation
$\rho_\tau=\mathbb{E}\langle\mathbf{r}_{c+\tau},\mathbf{r}_c\rangle/\sigma^2$.
Then
\[
    \mathbb{E}\|\mathbf{r}_{c+\tau}-\mathbf{r}_c\|^2
    =
    2\sigma^2(1-\rho_\tau),
\]
so the anchored residual term is smaller than the plain residual term
$\mathbb{E}\|\mathbf{r}_{c+\tau}\|^2=\sigma^2$ if and only if
$\rho_\tau>\tfrac{1}{2}$.
\end{proposition}

The consequence is practical. A model whose reconstruction error is large but
temporally persistent, so that it acts as a slowly varying bias, gains from
anchoring in proportion to that persistence. A model whose reconstruction error
behaves like white noise loses: at $\rho_\tau=0$ the anchored residual term is
twice the plain one in mean square, a factor $\sqrt{2}$ in norm. Since
$\rho_\tau$ is estimated directly from the residuals of a trained
autoencoder, without retraining or rolling anything out, the decision between
plain and anchored decoding can be made in advance for a given dataset and
horizon rather than discovered by trial.

% ============================================================
\section{Experimental Setup}
\label{sec:experimental}
% ============================================================

\subsection{Datasets and Protocol}

We evaluate the algorithm on two run-to-failure benchmarks from different
domains. The primary benchmark is NASA C-MAPSS
\citep{nasa_cmapss,saxena2008turbofan,saxena2008damage}, which consists of four
turbofan degradation subsets, FD001--FD004. Each engine trajectory starts from
an unknown initial wear state, operates normally at the beginning, develops a
fault, and is observed until failure in the training split. The test split is
truncated prior to failure and is accompanied by official RUL labels, which is
why the reconstruction, curvature, dimension, ablation and RUL studies are
carried out there.

The second benchmark is PHM~2010 high-speed CNC milling \citep{phm2010}, used in
Section~\ref{sec:milling} to test whether the two mechanisms of
Section~\ref{sec:theory} survive a change of domain. It replaces turbofan gas-path
measurements with force, vibration and acoustic-emission channels, and gradual
thermodynamic degradation with progressive tool wear, while retaining the
properties the algorithm requires. It carries no official RUL labels, so it
exercises stability and forecastability rather than downstream prognosis. The
criteria that make a dataset admissible are stated in
Section~\ref{sec:dataset-criteria}.

\begin{table}[t]
\centering
\caption{Characteristics of the C-MAPSS subsets used in this study.}
\label{tab:dataset_protocol}
\tablefont
\tighttables
\begin{tabular}{@{}lrrrrrr@{}}
\toprule
Dataset & Regimes & Faults & Train & Test & Sensors & Settings \\
\midrule
FD001 & 1 & 1 & 100 & 100 & 21 & 3 \\
FD002 & 6 & 1 & 260 & 259 & 21 & 3 \\
FD003 & 1 & 2 & 100 & 100 & 21 & 3 \\
FD004 & 6 & 2 & 249 & 248 & 21 & 3 \\
\bottomrule
\end{tabular}
\end{table}

The subsets form a useful complexity ladder. FD001 and FD003 have one
operating regime, while FD002 and FD004 have six regimes. FD003 and FD004 have
two fault modes. All subsets contain 21 sensor channels and three operating
settings.

\subsection{Models and Baselines}

We compare the proposed bounded latent manifold against sensor-space VAR and
VAR(2), GRU, LSTM, and TCN forecasting models, unbounded and partially
regularized autoencoder variants, higher-dimensional latent variants
($K=1,\ldots,6$), and specialist versus pooled generalized models. The baseline
comparison tests free-run stability, not merely one-step prediction accuracy.
VAR models are closed-loop sensor-space recursions. Neural sequence baselines
are rolled out recursively. The manifold model is rolled out in latent space
and decoded.

\subsection{Metrics}

We evaluate reconstruction, stability, forecastability, and downstream RUL.
Reconstruction is measured by $R^2$. Free-run stability is measured using
rollout growth ratios and a boundedness indicator. Forecastability is measured
by NRMSE at multiple horizons, skill relative to persistence, and latent
curvature $\kappa$. RUL performance is measured using RMSE, MAE, $R^2$, and
the NASA asymmetric score.

Final algorithm results use the production training budget. Comparative
ablation sweeps use a reduced epoch budget for tractability; their ordering and
qualitative conclusions match the production run. The hyperparameter
sensitivity study in Appendix~\ref{app:sensitivity} uses the full production
budget of 4000 autoencoder epochs over three seeds. All retraining scripts use
fixed seeds, with seed robustness evaluated over seeds $\{0,1,2,3,4\}$.

% ============================================================
\section{Results}
\label{sec:results}
% ============================================================

The results are organized around the algorithmic mechanism. We first establish
the learned latent geometry, then evaluate final performance, free-run
stability, dimension, ablations, robustness, transfer, and RUL context.

\subsection{Latent Manifold Quality}

A stable latent rollout is only meaningful if the representation organizes
degradation trajectories coherently. Fig.~\ref{fig:latent-manifold} shows the
intrinsic-dimensionality diagnostics and learned two-dimensional latent
manifolds for FD001--FD004.

\begin{figure}[t]
    \centering
    \includegraphics[width=\figlarge,height=0.43\textheight,keepaspectratio]{\mainfig/D2_manifold_FD001_FD004_combined.png}
    \caption{Intrinsic dimensionality and learned latent manifolds across C-MAPSS.}
    \label{fig:latent-manifold}
\end{figure}

The first two principal components explain 96.5\%, 79.5\%, 85.1\%, and 84.9\%
of variance on FD001--FD004, respectively. FD001 has the cleanest
two-dimensional structure, while FD002 and FD004 are more complex because of
multiple operating regimes. Nevertheless, the learned latent coordinates
organize trajectories by life fraction across all four subsets.

\subsection{Main Cross-Dataset Performance}

Table~\ref{tab:main-cross} summarizes the final bounded $K=2$ model.

\begin{table}[t]
\centering
\caption{Final model performance across C-MAPSS.}
\label{tab:main-cross}
\tinytable
\verytighttables
\begin{tabular}{@{}lrrrrrrrr@{}}
\toprule
Data & Reg. & $\rho_2$ & Rec. & $\rho_{\rm VAR}$ & VAR Gr. & Man. & RUL & RUL \\
     &      &          & $R^2$ &                 &          & Norm & RMSE & $R^2$ \\
\midrule
FD001 & 1 & .965 & .930 & 1.020 & 760 & 5.99 & 14.53 & .878 \\
FD002 & 6 & .795 & .865 & 1.016 & 468 & 2.29 & 27.02 & .748 \\
FD003 & 1 & .851 & .960 & 1.018 & 57.6 & 4.33 & 16.31 & .845 \\
FD004 & 6 & .849 & .930 & 1.015 & 25.0 & 3.58 & 27.58 & .744 \\
\bottomrule
\end{tabular}
\end{table}

Across all datasets, the model reconstructs selected degradation-sensitive
channels with mean reconstruction $R^2$ between 0.865 and 0.960. The fitted
sensor-space VAR models are non-contractive on every dataset, with
$\rho_{\mathrm{VAR}}>1$, and VAR free-run growth ranges from about $25\times$
to $760\times$. The bounded latent manifold remains controlled and supports
RUL prediction with RMSE from 14.53 to 27.58 cycles.

\subsection{Free-Run Stability Against Baselines}

The baseline study asks whether stable rollout is specific to the proposed
latent geometry or also appears in standard sensor-space and neural forecasters.

\begin{table}[t]
\centering
\caption{Free-run growth by model. Values are final/initial norm ratios.}
\label{tab:baseline-growth}
\tinytable
\verytighttables
\begin{tabular}{@{}lrrrrrr@{}}
\toprule
Data & Ours & VAR & VAR(2) & LSTM & GRU & TCN \\
\midrule
FD001 & 1.5 & 760 & 748 & 2.9 & 2.8 & 3.2 \\
FD002 & 1.5 & 468 & 476 & 4.7 & 13.0 & $1.27{\times}10^6$ \\
FD003 & 1.3 & 58 & 57 & 0.9 & 1.5 & 2.8 \\
FD004 & 0.8 & 25 & 27 & 3.4 & 1.9 & 1.4 \\
\bottomrule
\end{tabular}
\end{table}

Fig.~\ref{fig:baselines} summarizes the same comparison visually. The manifold
remains bounded on all four datasets. VAR and VAR(2) diverge on every dataset,
matching Theorem~\ref{thm:linear-growth}. Neural baselines are more flexible
but not uniformly stable: TCN diverges catastrophically on FD002, and GRU grows
by about $13\times$ on FD002.

\begin{figure}[t]
    \centering
    \includegraphics[width=\figmedium,height=0.34\textheight,keepaspectratio]{\mainfig/fig_baselines.png}
    \caption{Free-run growth comparison across rollout models.}
    \label{fig:baselines}
\end{figure}

At very short horizons, VAR can have competitive NRMSE because it is optimized
as a local linear predictor. The claim is therefore not strict one-step
superiority; it is stable long-horizon free-run behavior.

\subsection{Dataset Requirements for Cross-Domain Transfer}
\label{sec:dataset-criteria}
\label{app:dataset-criteria}

A candidate dataset for this algorithm must supply four things: operating
conditions or settings that are recorded separately from the measurements, so
that regimes can be identified and normalized out; noisy sensor channels, so
that denoising and trend-based selection have something to act on; a
degradation process, so that the health coordinate has a direction to track;
and enough temporal structure for the signal to be forecastable at the horizons
of interest. Official remaining-useful-life labels are \emph{not} required. RUL
is a downstream use of the latent state, not a precondition for learning it,
and a dataset that satisfies the first four properties can exercise
reconstruction, boundedness, curvature, and forecast skill without them.


\subsection{Cross-Domain Validation: CNC Tool Wear}
\label{sec:milling}

The results above are confined to turbofan degradation. To test whether the two
mechanisms of Section~\ref{sec:theory} survive a change of domain, we apply the
identical pipeline to the PHM~2010 CNC milling benchmark \citep{phm2010}: six
cutting tools run
to 315 cuts each, with thirty-five time-domain force, vibration and
acoustic-emission features per cut, four tools used for training and two held
out. The latent dimension is $K{=}4$ by the intrinsic-dimensionality elbow and
reconstruction $R^2$ is $0.730$; the fitted sensor-space companion matrix has
spectral radius $\rho=1.030$, so Theorem~\ref{thm:linear-growth} predicts
geometric divergence of the closed-loop sensor recursion.

\begin{table}[t]
\centering
\caption{CNC milling free-run growth. Values are final/initial decoded norm
ratios at two probe horizons: the data horizon (188 steps, the last observed
cut) and the end-of-life horizon (400 steps, matching the C-MAPSS protocol of
Table~\ref{tab:baseline-growth}). Rows marked \textbf{Ours} are the proposed
bounded latent model under the rollout rules of Section~\ref{sec:method}.}
\label{tab:milling-growth}
\tinytable
\verytighttables
\begin{tabular}{@{}llrcrc@{}}
\toprule
Model & & @188 & Bnd. & @400 & Bnd. \\
\midrule
persistence        & ref.            & 1.00 & yes & 1.00      & yes \\
sensor VAR         & ref.            & 2.69 & yes & $9104$    & \textbf{no} \\
\midrule
constant velocity  & \textbf{Ours} & 2.89 & yes & 1.62      & yes \\
damped CV          & \textbf{Ours} & 1.20 & yes & 0.98      & yes \\
Holt damped trend  & \textbf{Ours} & 1.52 & yes & 1.96      & yes \\
windowed trend     & \textbf{Ours} & 1.39 & yes & 1.00      & yes \\
anchored damped CV & \textbf{Ours} & 1.13 & yes & 0.98      & yes \\
anchored Holt      & \textbf{Ours} & 1.43 & yes & 2.14      & yes \\
\bottomrule
\end{tabular}
\end{table}

Table~\ref{tab:milling-growth} reproduces the central stability result outside
C-MAPSS, and it also shows why the probe horizon matters. Stopped at the last
observed cut, every model including the VAR appears bounded, and the ordering is
uninformative: plain constant-velocity rollout ($2.89$) looks slightly worse
than the VAR ($2.69$). Continued to the end-of-life horizon used throughout
Section~\ref{sec:results}, the ordering inverts decisively. The sensor-space
VAR grows by a factor of $9\,104$ and leaves the boundedness criterion, exactly
as Theorem~\ref{thm:linear-growth} requires of a non-contractive companion
matrix, while every bounded latent rollout remains controlled and two of them
contract. The distinction is not visible at short probe horizons because a
spectral radius of $1.030$ needs several hundred steps to express itself; it is
a property of the operator, not of the window in which it is observed.

\begin{table}[t]
\centering
\caption{CNC milling forecast skill relative to persistence. Positive values
beat persistence. Rows marked \textbf{Ours} are the proposed bounded latent
model; anchored rows use the decoding of Section~\ref{sec:anchored}.}
\label{tab:milling-skill}
\tinytable
\verytighttables
\begin{tabular}{@{}llrrrr@{}}
\toprule
Model & & @1 & @10 & @25 & @50 \\
\midrule
persistence        & ref.            & $0.00$    & $0.00$  & $0.00$  & $0.00$ \\
sensor VAR         & ref.            & $-9.51$   & $-3.71$ & $-2.63$ & $-2.86$ \\
\midrule
constant velocity  & \textbf{Ours} & $-263.2$  & $-9.16$ & $-5.06$ & $-0.87$ \\
damped CV          & \textbf{Ours} & $-263.2$  & $-6.38$ & $-2.43$ & $-0.35$ \\
Holt damped trend  & \textbf{Ours} & $-267.4$  & $-7.07$ & $-3.47$ & $-0.35$ \\
windowed trend     & \textbf{Ours} & $-263.4$  & $-7.42$ & $-3.81$ & $-0.74$ \\
anchored Holt      & \textbf{Ours} & $-1.35$   & $-1.74$ & $-1.78$ & $-0.12$ \\
anchored damped CV & \textbf{Ours} & $\mathbf{-0.73}$ & $\mathbf{-1.04}$ & $\mathbf{-0.86}$ & $\mathbf{+0.05}$ \\
\bottomrule
\end{tabular}
\end{table}

Table~\ref{tab:milling-skill} is the second mechanism, and it is the clearest
demonstration of Corollary~\ref{cor:anchor} in the paper. Every plain-decoded
rollout is catastrophic at horizon one, with skill near $-263$; in NRMSE the
error is $0.595$ against persistence at $0.0366$, a factor of sixteen. That
penalty is not caused by the latent dynamics, which at horizon one have barely
moved. It is the constant origin offset $\mathbf{r}_c=\y_c-D_\phi(\h_c)$ that
plain decoding carries at every horizon. Anchoring cancels it, and the same
rollout, decoded differently, moves from $-263.2$ to $-0.73$: an NRMSE of
$0.0482$, within a third of persistence rather than sixteen times worse.
Anchored damped constant velocity is also the only rule with positive skill at
the longest horizon.

Two features of this dataset explain why both mechanisms are visible on it at
once. The fitted damping is $\phi^\star=0.9$, close to undamped extrapolation,
so the latent rollout is genuinely doing work rather than contributing a single
frozen step; and the reconstruction residual is large enough that the offset
term dominates the curvature term at short horizons, which is the regime in
which Proposition~\ref{prop:anchor-gain} predicts anchoring to pay. The two
mechanisms of Section~\ref{sec:theory} are therefore both exercised on a single
external dataset: bounded latent geometry controls the admissible state space
over long free runs, and the decoding choice controls the error inside it.

\subsection{Latent Dimension: Capacity Versus Stability}

The $K$-aware dimension sweep tests whether the two-dimensional latent state is
a bottleneck or a dynamical regularizer. For each $K$, all latent coordinates
and velocities are used downstream, so higher-dimensional models are not
evaluated only through $(h_0,h_1)$.

\begin{figure}[t]
    \centering
    \includegraphics[width=\figmedium,height=0.34\textheight,keepaspectratio]{\mainfig/fig_k_tradeoff.png}
    \caption{Latent dimension tradeoff between reconstruction and prognosis.}
    \label{fig:k-tradeoff}
\end{figure}

The results support Proposition~\ref{prop:k-regularization}. Reconstruction
generally improves as $K$ increases, but stability and RUL do not improve
monotonically. On FD001, $K=2$ achieves horizon-50 NRMSE about 0.465, while
$K=4$ and $K=6$ rise to about 0.561 and 0.555. On FD002, larger latent spaces
become unbounded for $K\geq3$ under the stability criterion, even as
reconstruction improves. Thus the low-dimensional latent state acts as a
dynamical regularizer.

\subsection{Ablation: Boundedness Versus Forecastability}

The ablation study separates reconstruction, finite-horizon boundedness,
forecastability, and RUL utility.

\begin{table}[t]
\centering
\caption{Mechanism ablation on FD002.}
\label{tab:mechanism}
\tinytable
\verytighttables
\begin{tabular}{@{}lrrrr@{}}
\toprule
Variant & Rec. $R^2$ & $\kappa$ & Skill & RUL RMSE \\
\midrule
No regime norm. & .999 & -- & -0.13 & 40.5 \\
Unbounded AE & -- & .0346 & .276 & -- \\
Full model & .865 & .0019 & .676 & 27.02 \\
\bottomrule
\end{tabular}
\end{table}

\begin{figure}[t]
    \centering
    \includegraphics[width=\figmedium,height=0.34\textheight,keepaspectratio]{\mainfig/boundedness_mechanism.png}
    \caption{Mechanism analysis separating boundedness from forecastability.}
    \label{fig:boundedness-mechanism}
\end{figure}

The most important ablation result is negative but clarifying. Under local
constant-velocity latent rollout, the unbounded autoencoder remains finite
over the evaluated horizon because the rollout is open-loop with small
velocities and the tanh decoder saturates outside the training range. The
empirical advantage of the proposed model is therefore not that every
unbounded latent model diverges. The advantage is forecastability.

On FD002, the unbounded autoencoder has curvature $\kappa=0.0346$ and forecast
skill 0.276. The full model reduces curvature to 0.0019 and improves skill to
0.676. The no-regime-normalization variant reaches reconstruction
$R^2=0.999$ on FD002 but collapses to forecast skill $-0.13$ and RUL RMSE
40.5. This directly shows that reconstruction is not prognosis.

\subsection{Open-Loop Versus Closed-Loop Rollout}

\begin{table}[t]
\centering
\caption{Open-loop versus closed-loop rollout diagnostic.}
\label{tab:open-closed}
\tinytable
\verytighttables
\begin{tabular}{@{}llll@{}}
\toprule
Model & Space & Rollout & Behavior \\
\midrule
VAR/VAR(2) & sensor & closed-loop & diverges, $\rho>1$ \\
Unbounded AE & latent & open-loop CV & finite over horizon \\
Bounded AE & latent & projected CV & formally bounded \\
\bottomrule
\end{tabular}
\end{table}

This diagnostic clarifies why finite-horizon decoded boundedness can occur
even for the unbounded autoencoder. VAR and VAR(2) are recursive
closed-loop sensor-space models. The dramatic divergence in the baseline study
is associated with observation-space recursion and non-contractive transition
matrices, not with every unbounded latent parameterization.

\subsection{Robustness and Transfer}

The main conclusions persist across five random seeds. The full bounded model
remains bounded in 100\% of seed runs. FD001 RUL is $14.39\pm0.45$, and FD002
RUL is $27.40\pm0.13$. On FD002, the unbounded autoencoder has forecast skill
$-0.73\pm0.82$, while the full model achieves $0.52\pm0.13$.

\begin{table}[t]
\centering
\caption{Compact robustness and transfer summary.}
\label{tab:robust-transfer}
\tinytable
\verytighttables
\begin{tabular}{@{}lrrr@{}}
\toprule
Study & Dataset & Key metric & Value \\
\midrule
Seed robustness & FD001 & Full RUL RMSE & $14.39\pm0.45$ \\
Seed robustness & FD002 & Full RUL RMSE & $27.40\pm0.13$ \\
Seed robustness & FD002 & Unbounded AE skill & $-0.73\pm0.82$ \\
Seed robustness & FD002 & Full skill & $0.52\pm0.13$ \\
Transfer & FD002 & Specialist RUL RMSE & 26.87 \\
Transfer & FD002 & Generalized RUL RMSE & 26.92 \\
\bottomrule
\end{tabular}
\end{table}

The generalized model preserves bounded rollout on three of four datasets and
achieves near-identical FD002 RUL to the specialist model. This suggests that
the learned geometry is partially reusable across related degradation settings.
Section~\ref{sec:milling} tests the same pipeline outside the turbofan domain,
on CNC tool wear.

\subsection{Downstream RUL Utility and Literature Context}

Although the algorithm is not optimized solely as an RUL leaderboard model, the
learned latent state remains useful for downstream prognosis. The final model
reduces mean-baseline RMSE from 43.07--54.90 cycles to 14.53--27.58 cycles
across FD001--FD004.

\begin{figure}[t]
    \centering
    \includegraphics[width=\figlarge,height=0.38\textheight,keepaspectratio]{\mainfig/C3_health_vs_rul_FD001_FD004_combined.png}
    \caption{Learned latent health coordinate versus RUL.}
    \label{fig:health-rul}
\end{figure}

Table~\ref{tab:rul-literature} places the RUL numbers in context. External
numbers are reported by the cited papers and are not necessarily directly
comparable because preprocessing, RUL capping, validation choices, and target
definitions differ. The table is included only to show that the latent state
has useful downstream RUL signal; the central contribution remains stable
long-horizon rollout.

\begin{table}[t]
\centering
\caption{C-MAPSS RUL context. Values are RMSE in cycles.}
\label{tab:rul-literature}
\tinytable
\verytighttables
\begin{tabular}{@{}lrrrr@{}}
\toprule
Method & FD001 & FD002 & FD003 & FD004 \\
\midrule
Ours & 14.53 & 27.02 & 16.31 & 27.58 \\
AFICv \citep{alomari2023aficv} & 11.80 & 23.00 & 14.60 & 22.30 \\
CNN-LSTM-Att. \citep{deng2024cnnlstmattention} & 15.98 & 14.45 & 13.91 & 16.64 \\
TTSNet \citep{li2025ttsnet} & 11.02 & 13.25 & 11.06 & 18.26 \\
\bottomrule
\end{tabular}
\end{table}

% ============================================================
\section{Discussion}
\label{sec:discussion}
% ============================================================

\subsection{Boundedness Is Not Enough}

The ablation study shows that finite-horizon boundedness is not the complete
explanation for stable prognostic performance. Under constant-velocity latent
rollout, an unbounded autoencoder can remain numerically bounded over
practical horizons because the latent velocities are small and the tanh decoder
saturates. Thus, the empirical advantage of the proposed model is not merely
that its decoded trajectories remain finite.

The more important distinction is between boundedness and forecastability. The
bounded latent model provides a formal horizon-independent guarantee when the
full latent vector is projected into a compact domain. Smoothness and
monotonicity then reduce latent curvature, making constant-velocity forecasting
more accurate. Boundedness controls the admissible state space; smoothness
controls the predictability of motion within that state space.

\subsection{Reconstruction Is Not Prognosis}

Several experiments show that reconstruction accuracy can be misleading. The
no-regime-normalization ablation on FD002 achieves reconstruction $R^2=0.999$
but has negative forecast skill and poor RUL. Similarly, increasing latent
dimension improves reconstruction but can weaken rollout stability. This
supports the view that degradation representations should be evaluated by
forecastability and downstream prognostic utility, not reconstruction alone.

\subsection{Algorithmic Scope}

The proposed method is an algorithm for learning stable latent degradation
states from multivariate time series with operating variables. C-MAPSS is the
primary testbed because it provides controlled variation in operating regimes
and fault modes together with official labels, and CNC tool wear tests the same
mechanisms on a process with different sensing and different physics. The
algorithmic claim is not that the method is the best possible RUL regressor on
either. The claim is that bounded low-dimensional latent geometry and
low-curvature latent motion improve the reliability of long-horizon rollout,
and that where a decoded forecast is dominated by its origin reconstruction
offset, anchoring removes it; both claims are shown to hold across the two
domains, while useful downstream RUL information is preserved.

\subsection{Limitations}

The evaluation covers turbofan degradation (FD001--FD004) and CNC tool wear
(Section~\ref{sec:milling}). Downstream RUL is reported only on C-MAPSS, since
the milling benchmark carries no official remaining-useful-life labels, so the
cross-domain evidence covers stability and forecastability but not RUL
accuracy. Cross-domain evidence rests on a single external benchmark.

Anchored decoding is not universally preferable, and
Proposition~\ref{prop:anchor-gain} states the condition under which it helps.
Where the reconstruction residual is weakly correlated across the forecast
horizon the anchored bound is the looser of the two, so the choice between
plain and anchored decoding is a measurable property of a dataset rather than a
universal improvement.

The latent rollout model in the main results is also deliberately simple.
Constant-velocity extrapolation highlights the role of latent geometry, and the
learned bounded residual head of Section~\ref{sec:method} is exercised only in
the appendices; richer latent dynamics may improve accuracy further when
degradation is highly nonlinear or abrupt.

% ============================================================
\section{Conclusion}
% ============================================================

We presented a bounded latent manifold algorithm for stable multivariate
time-series forecasting in degradation settings. The method removes
operating-regime effects, denoises and filters degradation-sensitive sensors,
and learns a compact latent state with monotonicity and smoothness
regularization. Theoretical analysis separates sensor-space closed-loop
divergence, bounded latent rollout, and smoothness-controlled forecastability.
Experiments show that low-dimensional bounded manifolds act as dynamical
regularizers: larger latent dimensions can reconstruct better but roll out
worse, and boundedness alone is not sufficient without smooth latent motion.
The resulting representation provides stable long-horizon behavior and useful
downstream RUL prediction.

% ============================================================
\section*{Acknowledgements}
This section is intentionally left blank for submission and may be completed in
the camera-ready version if appropriate.

\bibliographystyle{cas-model2-names}
\bibliography{references}

% ============================================================
\appendix

\section{Full Experimental Tables}

The main text reports compact summaries to avoid page-wide tables. The full
machine-readable CSV files and LaTeX tables are included in the artifact bundle
under \texttt{paper\_artifacts/acml\_results\_bundle/}. They include the full
baseline table, full $K$-sweep, full ablation table, seed robustness table,
specialist-versus-generalized table, and implementation diagnostic table.

\section{Additional Rollout Diagnostics}

Fig.~\ref{fig:appendix-var-eigs} shows the VAR eigenvalue spectra used to
validate the non-contractive sensor-space rollout analysis. Fig.
\ref{fig:appendix-free-run} shows the corresponding free-run divergence traces.

\begin{figure}[t]
    \centering
    \includegraphics[width=\figappendix,height=0.48\textheight,keepaspectratio]{\appfig/A2_var_eigenvalues_FD001_FD004_combined.png}
    \caption{Sensor-space VAR eigenvalue spectra across FD001--FD004.}
    \label{fig:appendix-var-eigs}
\end{figure}

\begin{figure}[t]
    \centering
    \includegraphics[width=\figappendix,height=0.48\textheight,keepaspectratio]{\appfig/A4_free_run_divergence_FD001_FD004_combined.png}
    \caption{Free-run divergence trajectories for sensor-space VAR versus
    latent manifold rollout.}
    \label{fig:appendix-free-run}
\end{figure}

\begin{figure}[t]
    \centering
    \includegraphics[width=\figappendix,height=0.48\textheight,keepaspectratio]{\appfig/A1_rollout_r2_vs_horizon_FD001_FD004_combined.png}
    \caption{Rollout accuracy versus horizon.}
    \label{fig:appendix-rollout-horizon}
\end{figure}

\begin{figure}[t]
    \centering
    \includegraphics[width=\figappendix,height=0.48\textheight,keepaspectratio]{\appfig/A3_example_trajectories_FD001_FD004_combined.png}
    \caption{Representative sensor trajectories under true, manifold, and VAR
    rollouts.}
    \label{fig:appendix-example-trajectories}
\end{figure}

\section{Forecastability Diagnostics}

\begin{figure}[t]
    \centering
    \includegraphics[width=\figappendix,height=0.48\textheight,keepaspectratio]{\appfig/B1_health_forecasts_FD001_FD004_combined.png}
    \caption{Example latent health forecasts.}
    \label{fig:appendix-health-forecasts}
\end{figure}

\begin{figure}[t]
    \centering
    \includegraphics[width=\figappendix,height=0.48\textheight,keepaspectratio]{\appfig/B2_error_vs_horizon_FD001_FD004_combined.png}
    \caption{Forecast error growth versus horizon.}
    \label{fig:appendix-error-horizon}
\end{figure}

\begin{figure}[t]
    \centering
    \includegraphics[width=\figappendix,height=0.48\textheight,keepaspectratio]{\appfig/B3_skill_vs_persistence_FD001_FD004_combined.png}
    \caption{Forecast skill relative to persistence.}
    \label{fig:appendix-skill}
\end{figure}

\section{RUL Diagnostics}

\begin{figure}[t]
    \centering
    \includegraphics[width=\figappendix,height=0.48\textheight,keepaspectratio]{\appfig/C1_rul_scatter_FD001_FD004_combined.png}
    \caption{Predicted versus true RUL and error distributions.}
    \label{fig:appendix-rul-scatter}
\end{figure}

\begin{figure}[t]
    \centering
    \includegraphics[width=\figappendix,height=0.48\textheight,keepaspectratio]{\appfig/C2_examples_FD001_FD004_combined.png}
    \caption{Representative test-engine RUL examples.}
    \label{fig:appendix-rul-examples}
\end{figure}

\section{Additional Latent Geometry}

\begin{figure}[t]
    \centering
    \includegraphics[width=\figappendix,height=0.48\textheight,keepaspectratio]{\appfig/D1_health_trajectories_FD001_FD004_combined.png}
    \caption{Primary health-coordinate trajectories across engines.}
    \label{fig:appendix-health-trajectories}
\end{figure}

\begin{figure}[t]
    \centering
    \includegraphics[width=\figappendix,height=0.48\textheight,keepaspectratio]{\appfig/k_vs_reconstruction_and_stability.png}
    \caption{Expanded latent-dimension sweep showing reconstruction and
    stability axes.}
    \label{fig:appendix-k-stability}
\end{figure}

\begin{figure}[t]
    \centering
    \includegraphics[width=\figappendix,height=0.48\textheight,keepaspectratio]{\appfig/k_vs_rul.png}
    \caption{RUL RMSE versus latent dimension.}
    \label{fig:appendix-k-rul}
\end{figure}

\section{Hyperparameter Sensitivity Analysis}
\label{app:sensitivity}

This appendix reports a full-budget sensitivity analysis for the four core
hyperparameters of the algorithm: the monotonicity penalty $\lambda_m$, the
smoothness penalty $\lambda_s$, the latent-rollout step size $\gamma$ (the
residual dynamics coefficient $\alpha$), and the number of normalization
regimes. Each parameter is swept one at a time around the production defaults
($\lambda_m{=}5$, $\lambda_s{=}2$, $\gamma{=}0.05$, regimes selected
automatically), with every configuration trained at the full production budget
of 4000 autoencoder epochs and averaged over three seeds. Representative
single-regime (FD001) and multi-regime (FD002) subsets are used, so the trends
below are obtained under the same schedule as the final model.
Fig.~\ref{fig:appendix-sensitivity} summarizes all four sweeps.

\begin{figure}[t]
    \centering
    \includegraphics[width=\figwide,height=0.30\textheight,keepaspectratio]{\appfig/sensitivity_summary.png}
    \caption{Full-budget hyperparameter sensitivity (mean $\pm$ std over three
    seeds). (a) monotonicity penalty, (b) smoothness penalty, (c) latent step
    size $\gamma$, (d) regime count.}
    \label{fig:appendix-sensitivity}
\end{figure}

\begin{table}[t]
\centering
\caption{Sensitivity to the monotonicity penalty $\lambda_m$ (with
$\lambda_s{=}2$). Full budget, mean over three seeds. Skill@20 is
constant-velocity forecast skill of the health coordinate relative to
persistence; $\kappa$ is latent curvature (lower is smoother).}
\label{tab:sens-mono}
\tinytable
\verytighttables
\begin{tabular}{@{}llrrrr@{}}
\toprule
Data & $\lambda_m$ & Recon $R^2$ & $\kappa$ & Skill@20 & RUL RMSE \\
\midrule
FD001 & 0  & 0.969 & $2.9\!\times\!10^{-3}$ & $+0.735$ & 14.94 \\
FD001 & 1  & 0.969 & $1.5\!\times\!10^{-3}$ & $+0.749$ & 14.67 \\
FD001 & 5  & 0.969 & $6.2\!\times\!10^{-4}$ & $+0.751$ & 14.47 \\
FD001 & 10 & 0.968 & $3.5\!\times\!10^{-4}$ & $+0.769$ & 14.62 \\
FD002 & 0  & 0.919 & $5.8\!\times\!10^{-3}$ & $-0.092$ & 28.04 \\
FD002 & 1  & 0.918 & $2.4\!\times\!10^{-3}$ & $+0.286$ & 27.63 \\
FD002 & 5  & 0.914 & $1.5\!\times\!10^{-3}$ & $+0.652$ & 27.08 \\
FD002 & 10 & 0.911 & $8.5\!\times\!10^{-4}$ & $+0.550$ & 27.28 \\
\bottomrule
\end{tabular}
\end{table}

The monotonicity penalty is the most consequential of the two autoencoder
penalties. On the single-regime FD001 subset reconstruction is essentially
flat across $\lambda_m$ while curvature falls monotonically, so larger values
trade nothing measurable for a smoother health coordinate. On the multi-regime
FD002 subset the effect is sharper: with $\lambda_m{=}0$ the health coordinate
is not constrained to progress with life and forecast skill collapses (mean
$-0.09$, with high seed variance), whereas the production value $\lambda_m{=}5$
recovers stable positive skill ($+0.65$). Very large penalties ($\lambda_m{=}10$)
begin to erode reconstruction. The default sits at the knee of this tradeoff.

\begin{table}[t]
\centering
\caption{Sensitivity to the smoothness penalty $\lambda_s$ (with
$\lambda_m{=}5$). Full budget, mean over three seeds.}
\label{tab:sens-smooth}
\tinytable
\verytighttables
\begin{tabular}{@{}llrrrr@{}}
\toprule
Data & $\lambda_s$ & Recon $R^2$ & $\kappa$ & Skill@20 & RUL RMSE \\
\midrule
FD001 & 0   & 0.968 & $6.2\!\times\!10^{-4}$ & $+0.752$ & 14.16 \\
FD001 & 0.5 & 0.969 & $6.1\!\times\!10^{-4}$ & $+0.752$ & 14.22 \\
FD001 & 2   & 0.969 & $6.2\!\times\!10^{-4}$ & $+0.751$ & 14.47 \\
FD001 & 5   & 0.969 & $6.4\!\times\!10^{-4}$ & $+0.754$ & 14.22 \\
FD002 & 0   & 0.915 & $1.6\!\times\!10^{-3}$ & $+0.648$ & 27.10 \\
FD002 & 0.5 & 0.906 & $1.5\!\times\!10^{-3}$ & $+0.649$ & 27.16 \\
FD002 & 2   & 0.914 & $1.5\!\times\!10^{-3}$ & $+0.652$ & 27.08 \\
FD002 & 5   & 0.915 & $1.5\!\times\!10^{-3}$ & $+0.654$ & 27.02 \\
\bottomrule
\end{tabular}
\end{table}

The smoothness penalty is a low-sensitivity knob. Across the full range on both
subsets, reconstruction, curvature, forecast skill, and RUL are all essentially
unchanged. This indicates that once the monotonicity constraint fixes the
orientation of the health coordinate, the additional roughness penalty is a
mild regularizer rather than a delicate tuning parameter.

\begin{table}[t]
\centering
\caption{Sensitivity to the latent step size $\gamma$ ($\alpha$). Full budget,
mean over three seeds. Entries are learned-dynamics rollout skill of the health
coordinate relative to persistence at horizons 10, 25, and 50.}
\label{tab:sens-gamma}
\tinytable
\verytighttables
\begin{tabular}{@{}llrrr@{}}
\toprule
Data & $\gamma$ & Skill@10 & Skill@25 & Skill@50 \\
\midrule
FD001 & 0.02 & $+0.41$ & $+0.59$ & $+0.63$ \\
FD001 & 0.05 & $+0.24$ & $+0.40$ & $+0.47$ \\
FD001 & 0.10 & $+0.43$ & $+0.63$ & $+0.64$ \\
FD001 & 0.25 & $-8.99$ & $-1.28$ & $+0.06$ \\
FD002 & 0.02 & $+0.29$ & $+0.62$ & $+0.61$ \\
FD002 & 0.05 & $+0.33$ & $+0.67$ & $+0.68$ \\
FD002 & 0.10 & $+0.33$ & $+0.66$ & $+0.68$ \\
FD002 & 0.25 & $-5.39$ & $-0.53$ & $+0.19$ \\
\bottomrule
\end{tabular}
\end{table}

The step size $\gamma$ separates a stable band from an unstable one rather than
identifying a sharp optimum. Every value up to $0.10$ yields positive skill at
all three horizons on both subsets, while the largest value tested,
$\gamma{=}0.25$, destabilizes the recursive rollout and produces strongly
negative short-horizon skill. Within the stable band the settings are not
separated by three seeds: on FD001 the spread at $\gamma{=}0.05$ (skill@10
$0.24\pm0.38$) overlaps both $\gamma{=}0.02$ ($0.41\pm0.09$) and
$\gamma{=}0.10$ ($0.43\pm0.12$), and on FD002 the same three settings are
indistinguishable. The step size therefore needs to be kept below the stability
threshold and does not otherwise require tuning.

\begin{table}[t]
\centering
\caption{Sensitivity to the number of normalization regimes. Full budget, mean
over three seeds. FD001 is physically single-regime; FD002 has six operating
conditions. ``Bnd.'' is the fraction of the three seeds whose decoded free-run
growth stays under the boundedness threshold.}
\label{tab:sens-regimes}
\tinytable
\verytighttables
\begin{tabular}{@{}llrrrrr@{}}
\toprule
Data & Regimes & Recon $R^2$ & Skill@20 & RUL RMSE & Growth & Bnd. \\
\midrule
FD001 & 1 & 0.969 & $+0.751$ & 14.47 & 1.16 & 3/3 \\
FD001 & 2 & 0.967 & $+0.754$ & 14.21 & 1.16 & 3/3 \\
FD001 & 3 & 0.969 & $+0.751$ & 14.26 & 1.25 & 3/3 \\
FD001 & 6 & 0.969 & $+0.754$ & 14.14 & 1.43 & 3/3 \\
FD002 & 1 & 0.986 & $-0.294$ & 44.64 & 1.83 & 3/3 \\
FD002 & 2 & 0.962 & $-13.67$ & 40.27 & 3.39 & 2/3 \\
FD002 & 3 & 0.901 & $-0.858$ & 32.00 & 12.69 & 0/3 \\
FD002 & 6 & 0.914 & $+0.652$ & 27.08 & 2.04 & 3/3 \\
\bottomrule
\end{tabular}
\end{table}

The regime count is the most safety-critical hyperparameter, and it is also the
clearest illustration of the paper's central distinction between reconstruction
and prognosis. On the single-regime FD001 subset the number of regimes is
irrelevant: forcing up to six clusters on data with one operating condition
leaves every metric unchanged. On the six-condition FD002 subset the behaviour
is the opposite and is initially counter-intuitive. Using too few regimes gives
\emph{higher} reconstruction ($R^2{=}0.986$ at one regime) but negative
forecast skill and RUL error above 44 cycles, because the autoencoder spends
its capacity reproducing the operating-condition staircase rather than the
degradation signal. Only when the regime count matches the six physical
operating conditions does forecast skill become positive ($+0.65$) and RUL
error drop to 27 cycles, despite a lower reconstruction $R^2$. High
reconstruction with too few regimes is therefore actively misleading.

The growth column shows the same effect on the stability side. At three regimes
on FD002 the decoded free-run growth reaches $12.7$ and falls outside the
boundedness flag in all three seeds, and at two regimes in one seed of three.
This does not violate Theorem~\ref{thm:bounded-rollout}, which concerns the
image $D_\phi(\Hspace)$ of the compact latent box and holds whatever the regime
count; the flag is the stricter empirical criterion used throughout, namely
whether the decoded trajectory norm grows by less than a factor of five over a
400-step free run. A mismatched regime count leaves residual
operating-condition structure in the normalized signal, the health coordinate
is driven against the faces of the box, and the decoded trajectory traverses a
large part of a set that is nonetheless compact. The architectural guarantee
survives a badly chosen regime count; the useful behaviour does not.

\section{Implementation Diagnostics}

The theory-matched full-box latent projection and the legacy $h_0$-only
projection are compared in the implementation diagnostic artifact
\texttt{theory\_impl\_rollout\_compare}, which shows materially similar rollout
metrics for the two projection rules.


\end{document}
